% NIH Specific Aims writing scaffold; not an agency-issued form.
% Reviewed: 2026-10-01. This compact draft compiles to one page with the supplied
% placeholders; actual content must be counted again. The NOFO controls.
% https://grants.nih.gov/grants-process/write-application/how-to-apply-application-guide/page-limits
% https://grants.nih.gov/grants-process/write-application/how-to-apply-application-guide/format-attachments
\documentclass[11pt,letterpaper]{article}
\usepackage[margin=0.5in]{geometry}
\usepackage{helvet}
\renewcommand{\familydefault}{\sfdefault}
\pagestyle{empty}
\setlength{\parindent}{0pt}
\setlength{\parskip}{8pt}

\begin{document}
\begin{center}
\textbf{Your Project Title Here: Concise and Descriptive}
\end{center}

\textbf{Significance and gap.} [Describe the important problem and the specific
unresolved issue using verified evidence. Explain why the proposed work is needed
without inventing disease-burden statistics or promising an established cure.]

\textbf{Objective and premise.} [State the objective of this proposal and the
hypothesis or other guiding premise appropriate to the mechanism. Explain the
rationale and feasibility using actual prior evidence or preliminary work where
applicable. Exploratory mechanisms do not universally require preliminary data.]

\textbf{Aim 1: [State a specific objective].}
[Describe the approach, the evidence it will produce, and how that evidence will
address the objective. State an expected outcome as a prediction, not a completed
result. Identify the principal alternative or decision criterion.]

\textbf{Aim 2: [State a related, testable objective].}
[Describe the approach and expected contribution. Explain how this aim remains
informative if Aim 1 has an unexpected result, where relevant.]

% Add or remove aims to fit the mechanism and scientific scope, not a fixed quota.
% \textbf{Aim 3: [Optional objective].} [Approach and expected contribution.]

\textbf{Expected impact.} [Explain what the completed work would establish, what
would remain unresolved, and how the result enables the next scientific step.]

\end{document}

% Current general NIH attachment requirements:
% - Specific Aims: generally one page, unless the NOFO specifies otherwise.
% - Text at least 11 pt; at most 15 characters per linear inch and six lines per
%   vertical inch. Smaller text in graphics is allowed if legible at 100%.
% - At least 0.5 inch margins; no added header/footer/page number unless requested.
% - Confirm the rendered PDF; a font family/package name does not prove compliance.
% - Do not shrink fonts or spacing to make an overlong draft fit.
